\documentclass[11pt]{article}
\usepackage[margin=1in]{geometry}
\usepackage{amsmath,amssymb,amsthm,mathtools,booktabs,hyperref}
\newtheorem{definition}{Definition}
\title{Phosphorus, Agriculture, and Catchment Safety:\\A Multi-Scale Material--Service--Institution Module}
\author{Author(s) redacted for review}
\date{}
\begin{document}\maketitle

\begin{abstract}
This paper specifies a multi-scale phosphorus module for the hybrid sustainability-systems architecture. It separates extraction and processing, trade, regional agriculture and consumption, recovery and waste, and catchment receiving-water dynamics. Elemental phosphorus is tracked as a declared material moiety across all represented compartments, while soil function and aquatic ecological condition are modeled as functional states rather than silently inferred from crop yield or total phosphorus. The module treats trade, transport, erosion, recovery, incomplete observations, institutional authority, and distributional constraints as explicit interfaces. It is a falsifiable template: no claim is made that a single regional model determines global reserve depletion or that agricultural output certifies soil or water safety.
\end{abstract}

\section{Scope and spatial boundary}
A phosphorus system cannot place global rock reserves, regional farms, and a local eutrophic lake in one unstructured lumped stock. The module therefore has three linked scales:
\[
R^{\rm ext}\xrightarrow{T^{\rm trade}}R^{\rm region}\xrightarrow{T^{\rm land}}R^{\rm catch}.
\]
The extraction layer contains mines, processing, and export nodes; the regional layer contains fertilizer, crop, livestock, food/product, waste, and recovery compartments; the catchment layer contains soil, erosion, river/lake, and ecological-function states. Each layer declares its boundary, time resolution, source data, and interface flows.

\section{Material state and stoichiometric routing}
Choose one elemental-phosphorus mass unit (for example kg P) and one time unit before assembling flows; conversion from phosphate or wet-product units uses explicit stoichiometric and moisture coefficients. All represented phosphorus compartments are in a single material vector
\[
r_P=(P_{\rm reserve},P_{\rm process},P_{\rm fertilizer},P_{\rm soil,labile},P_{\rm soil,stable},P_{\rm crop},P_{\rm livestock},P_{\rm food/product},P_{\rm waste},P_{\rm recovery},P_{\rm river/lake},\ldots)^\top.
\]
The exact components depend on the selected region and catchment. The material equation is
\[
\dot r_P=\mathsf S_P\nu_P(X_t,U,\omega,\vartheta)+b_P(X_t,U,\omega,t),
\]
with elemental-phosphorus moiety row $\ell_P^\top\mathsf S_P=0$. The boundary term is decomposed as $b_P=b_P^{\rm meas}+b_P^{\rm est}+e_P^{\rm report}+\delta_P^{\rm struct}$, separating measured and estimated physical flows from reporting error and structural discrepancy. Only physical boundary terms enter the material balance as mass transfer. Trade, imports, exports, and atmospheric/dust pathways appear as signed physical boundary flows. Every yield loss, product transformation, erosion pathway, recovery process, and event jump is routed explicitly.

At an event time $t_k$, the hybrid balance is
\[
\ell_P^\top[r_P(t_k^+)-r_P(t_k^-)]=J_{P,k}^{\rm boundary},
\]
where an internal reclassification has $J_{P,k}^{\rm boundary}=0$. Summing flow and jump terms gives the total phosphorus identity over any interval. Reporting corrections alter the estimate or belief, not physical mass, unless a physical boundary transfer is separately documented.

Reverse transformations have separate nonnegative elementary fluxes. Physical withdrawal and removal pathways are donor-limited. Every event reset maps the nonnegative compartment domain into itself and respects donor availability and recipient capacity; otherwise it is inadmissible. Material closure does not assert energetic feasibility of mining, fertilizer processing, transport, recovery, or recycling.

\section{Functional states and service maps}
Let
\[
f_P=(f_{\rm soil},f_{\rm aquatic},f_{\rm habitat})
\]
contain soil fertility/structure, aquatic ecological condition, and habitat effects. Any additional functional coordinate must have a definition, direction, dynamics, and measurement relation before it is added. These are not inferred automatically from material totals.

Agricultural, nutritional, and water-quality services are typed outputs:
\[
s=\mathcal F_P(r_P,f_P,U,\omega,\vartheta).
\]
Crop yield may be one service output but is not a certificate of soil-function, micronutrient, biodiversity, or receiving-water safety. A claimed substitute or recovery technology must appear as a material, energetic, temporal, waste, trade, and institutional pathway; no aggregate elasticity is treated as proof of substitution.

\section{Safety, service, and distributional constraints}
State safety may include, for example,
\[
\mathcal C_P(\lambda)=\{(r_P,f_P):P_{\rm soil}^{\min}\le P_{\rm soil,labile}\le P_{\rm soil}^{\max},\ f_{\rm soil}\ge f_{\rm soil}^{\min},\ f_{\rm aquatic}\ge f_{\rm aquatic}^{\min},\ P_{\rm river/lake}\le P_{\rm water}^{\max},\ldots\}.
\]
A service requirement and allocation constraint belong in the state-control relation
\[
\mathcal A_P(X,\omega;\lambda)=\{U:s(X,U,\omega)\ge s^{\min}(\lambda),\ \mathcal D_{\rm alloc}(X,U)\in\mathfrak D_{\rm adm}(\lambda)\}.
\]
Here $\mathcal D_{\rm alloc}$ records declared rights, burden, affordability, or procedural outcomes and is not an action correspondence. The module does not derive $\lambda$ mathematically; it exposes the viability consequences of stated normative choices.

\section{Trade, routing, and compositional interfaces}
Trade and transport are matrices/correspondences, not implicit background processes. Let $T^{\rm trade}_{ij}$ map extraction/processing node $i$ to region $j$, and let $T^{\rm land}$ route regional inputs, products, waste, and erosion to catchments. Interfaces must specify sign, units, delay, ownership, and boundary status.

A regional safety claim must check displacement: a local reduction in water pollution or fertilizer use can export phosphorus burden through imports, trade, waste shipment, or feed supply. Compositional certificates require interface assumptions, shared-control compatibility, and nonblocking institutional events.

\section{Observation, compatible states, and hybrid institutions}
Observations may include mine/production statistics, trade records, fertilizer sales, farm surveys, soil tests, crop sampling, waste/recovery data, river/lake chemistry, remote sensing, and ecological monitoring. Their relation to latent state is
\[
Y_k=\mathcal O_P(X_{[t_{k-1},t_k]})+\varepsilon_k.
\]
A compatible-state set
\[
B_k\subseteq C([t_{k-1}-\tau,t_k],\mathcal X_P)\times\Theta
\]
is a nonempty compact subset of a declared history/parameter metric space and contains all histories and parameters consistent with declared records, errors, controls, implementation branches, disturbances, and structural discrepancy. Its set-valued update has a closed graph; if compactness or closure is not established, fixed-point kernel claims are not made. Soil organic carbon, crop yield, nutrient concentration, or satellite imagery are proxies only after a measurement relation and validation error are specified.

Institutional state records authority over mining, trade, fertilizer regulation, nutrient management, recovery infrastructure, water-quality enforcement, monitoring, and allocation. At review events a prescription is selected under authority constraints and may yield multiple implementation/compliance outcomes. Instantaneous actions belong to an admissible correspondence, while a causal policy maps available observation histories to prescriptions; the two are not conflated.

\section{Competing hypotheses and identification}
The minimum competing model set should include: (H0) an aggregate regional phosphorus balance; (H1) a multi-compartment regional/catchment model; and (H2) a spatial trade-network/catchment model. Before fitting, declare held-out log score or proper scoring rule, coverage/calibration targets, threshold-event sensitivity/specificity where events exist, and a noninferiority or superiority margin. The more complex model is supported only by preregistered improvement, calibrated uncertainty, safety-boundary relevance, or a demonstrated decision advantage.

Identification must distinguish process uncertainty, parameter uncertainty, observation error, reporting error, structural discrepancy, trade leakage, and implementation uncertainty. A fitted residual may not be silently assigned to recovery, erosion, unreported trade, or soil immobilization.

\section{Validation and falsification}
A preregistration must declare material boundary, moiety units, stock/flow definitions, interface matrices, safety thresholds, observation/proxy model, uncertainty semantics, institutional events, comparator models, held-out targets, and falsification criteria.

For the preregistered task, boundary, resolution, and decision criterion, the module is rejected as a necessary architecture if simpler models predict held-out material/service/safety outcomes equally well and provide equal or better calibrated uncertainty. Conversely, descriptive evidence of nutrient depletion, eutrophication, or yield change does not itself validate the full framework unless the model-defined states and constraints are measured.

\section{Diagnostics, recovery, and safe learning}
Diagnostics distinguish gross material throughput, net stock decline, modeled threshold-crossing time, and pathwise failure risk. A local recovery intervention must be assessed for catchment and trade displacement. Recovery requires return to a viable material-functional-information-institution state, not merely temporary yield restoration.

Monitoring or pilot recovery actions may be safely informative only if they remain inside the declared safety envelope while reducing compatible-state uncertainty. This is a future adaptive-management test, not an assumption that learning is harmless.

\section{Conclusion}
This module preserves phosphorus as a material moiety while allowing soil function, aquatic condition, trade, services, authority, and implementation to matter independently. It offers a deliberately hard test of the general architecture: if the typed multi-scale model does not improve accounting, observability, or safety decisions relative to credible simpler alternatives, the architecture must be narrowed.

\end{document}
